% =========================================================
% Algebraic Structures — Relations
% =========================================================
% This file treats relations as algebraic objects carrying
% behaviours (properties, closure conditions, composition).
% Set-theoretic foundations (ordered pairs, Cartesian product,
% subset) are assumed from Volume I, Sets–Relations–Functions.
% The focus here is on what relations *do*, not what they *are*.

\section{Relations as Algebraic Objects}

% ---------------------------------------------------------
% TOOLKIT
% ---------------------------------------------------------
\begin{toolkitbox}{Relations — Quick Reference}
\small
\begin{tabular}{l l l}
\toprule
\textbf{Concept / Property} & \textbf{Definition / Condition} & \textbf{Detail} \\
\midrule
Unary relation     & Predicate $P$ on elements of $A$; subset $S \subseteq A$            & \hyperref[def:unary-relation]{↓} \\
$n$-ary relation   & Predicate on $n$-tuples; subset of $A^n$                            & \hyperref[def:nary-relation]{↓} \\
Binary relation    & Predicate on pairs; $R \subseteq A \times A$                        & \hyperref[def:binary-relation]{↓} \\
\midrule
Reflexive      & $\forall a \in A,\quad a \mathrel{R} a$                                  & \hyperref[def:rel-reflexive]{↓} \\
Irreflexive    & $\forall a \in A,\quad \lnot(a \mathrel{R} a)$                           & \hyperref[def:rel-irreflexive]{↓} \\
Symmetric      & $a \mathrel{R} b \Rightarrow b \mathrel{R} a$                             & \hyperref[def:rel-symmetric]{↓} \\
Antisymmetric  & $a \mathrel{R} b \land b \mathrel{R} a \Rightarrow a = b$                 & \hyperref[def:rel-antisymmetric]{↓} \\
Asymmetric     & $a \mathrel{R} b \Rightarrow \lnot(b \mathrel{R} a)$                      & \hyperref[def:rel-asymmetric]{↓} \\
Transitive     & $a \mathrel{R} b \land b \mathrel{R} c \Rightarrow a \mathrel{R} c$       & \hyperref[def:rel-transitive]{↓} \\
Total          & $\forall a,b \in A,\quad a \mathrel{R} b \lor b \mathrel{R} a$            & \hyperref[def:rel-total]{↓} \\
\midrule
Composition    & $(S \circ R)(a,c) \iff \exists b,\ a\mathrel{R}b \land b\mathrel{S}c$    & \hyperref[def:rel-composition]{↓} \\
\midrule
Equivalence relation & Reflexive + Symmetric + Transitive                                & \hyperref[def:equiv-relation]{↓} \\
Partial order        & Reflexive + Antisymmetric + Transitive                            & \hyperref[def:partial-order]{↓} \\
Strict partial order & Irreflexive + Asymmetric + Transitive                             & \hyperref[def:strict-partial-order]{↓} \\
Total order          & Partial order + Total                                             & \hyperref[def:total-order]{↓} \\
\bottomrule
\end{tabular}
\end{toolkitbox}


% ---------------------------------------------------------
% Unary relation
% ---------------------------------------------------------

\begin{definitionbox}{Definition (Unary Relation on a Set)}
\label{def:unary-relation}
Let $A$ be a set.
A \emph{unary relation on $A$} is a predicate $P$ that determines,
for each element $a \in A$, whether $a$ has the property $P$.
We write $P(a)$ when the property holds.
\end{definitionbox}

\begin{remark}[Fully quantified form]
$P$ is a unary relation on $A$
$\;\iff\;$
$(\forall a \in A)\;[\,P(a) \text{ is either true or false}\,]$.
\end{remark}

\begin{remark}[Unary relations are subsets]
A unary relation on $A$ corresponds exactly to a subset $S \subseteq A$:
the elements for which $P$ holds.
Equivalently, $P$ is the characteristic predicate of $S$.
Examples: ``is even'' on $\mathbb{Z}$; ``is invertible'' in a ring;
``is prime'' on $\mathbb{N}$.
Unary relations carry no properties like associativity or commutativity —
those belong to operations.
\end{remark}

% ---------------------------------------------------------
% n-ary relation
% ---------------------------------------------------------

\begin{definitionbox}{Definition ($n$-ary Relation on a Set)}
\label{def:nary-relation}
Let $A$ be a set and $n \in \mathbb{N}$.
An \emph{$n$-ary relation on $A$} is a predicate $R$ on $n$-tuples from $A$:
it determines, for each $(a_1, \ldots, a_n) \in A^n$, whether the tuple
satisfies $R$.
We write $R(a_1, \ldots, a_n)$ when it does.
\end{definitionbox}

\begin{remark}[Fully quantified form]
$R$ is an $n$-ary relation on $A$
$\;\iff\;$
$(\forall (a_1, \ldots, a_n) \in A^n)\;
[\,R(a_1, \ldots, a_n) \text{ is either true or false}\,]$.
\end{remark}

\begin{remark}[Special cases]
A unary relation ($n = 1$) is a property of individual elements.
A binary relation ($n = 2$) is a relationship between pairs — the case
carrying the richest algebraic structure (reflexivity, symmetry, transitivity, etc.).
A ternary relation ($n = 3$) arises in geometry (betweenness), algebra
(the relation $a + b = c$ on $\mathbb{Z}$), and logic.
The $n$-ary framework is the common abstraction.
\end{remark}

% ---------------------------------------------------------
% Binary relation on a set (algebraic restatement)
% ---------------------------------------------------------

\begin{definitionbox}{Definition (Binary Relation on a Set)}
\label{def:binary-relation}
Let $A$ be a set.
A \emph{binary relation on $A$} is a rule $R$ that determines,
for each ordered pair $(a, b)$ with $a, b \in A$, whether
$a$ and $b$ stand in the relation.
We write $a \mathrel{R} b$ when the relation holds, and
$\lnot(a \mathrel{R} b)$ when it does not.
\end{definitionbox}

\begin{remark}[Fully quantified form]
$R$ is a binary relation on $A$
$\;\iff\;$
$(\forall a \in A)(\forall b \in A)\;
[\,a \mathrel{R} b \;\text{ is either true or false}\,]$.
\end{remark}

\begin{remark}[Set-theoretic grounding]
In Volume I we identified a binary relation on $A$ with a subset
$R \subseteq A \times A$, writing $(a,b) \in R$ in place of
$a \mathrel{R} b$.
The infix notation $a \mathrel{R} b$ adopted here suppresses that
representation and treats the relation as a primitive predicate on
pairs — the algebraic point of view.
Familiar examples: $=$, $\leq$, $\mid$ (divides), $\subseteq$.
\end{remark}

% ---------------------------------------------------------
% Element-level properties
% ---------------------------------------------------------

\begin{definitionbox}{Definition (Reflexive / Irreflexive)}
\label{def:rel-reflexive}
\label{def:rel-irreflexive}
Let $R$ be a binary relation on $A$.
\begin{enumerate}[label=(\roman*)]
  \item $R$ is \emph{reflexive} if every element of $A$ is related to itself:
        \[
          \forall a \in A,\quad a \mathrel{R} a.
        \]
  \item $R$ is \emph{irreflexive} if no element of $A$ is related to itself:
        \[
          \forall a \in A,\quad \lnot(a \mathrel{R} a).
        \]
\end{enumerate}
\end{definitionbox}

\begin{remark}[Fully quantified form]
Reflexive: $(\forall a \in A)\; a \mathrel{R} a$.
\quad
Irreflexive: $(\forall a \in A)\; \lnot(a \mathrel{R} a)$.
\end{remark}

\begin{remark}[Independence]
Reflexive and irreflexive are \emph{not} negations of each other.
A relation on a nonempty set cannot be both reflexive and irreflexive.
A relation on the empty set is vacuously both.
The relation $\{(1,2)\}$ on $\{1,2\}$ is neither.
\end{remark}

\begin{definitionbox}{Definition (Symmetric / Antisymmetric / Asymmetric)}
\label{def:rel-symmetric}
\label{def:rel-antisymmetric}
\label{def:rel-asymmetric}
Let $R$ be a binary relation on $A$.
\begin{enumerate}[label=(\roman*)]
  \item $R$ is \emph{symmetric} if relatedness in one direction implies
        relatedness in both:
        \[
          \forall a,b \in A,\quad a \mathrel{R} b \;\Rightarrow\; b \mathrel{R} a.
        \]
  \item $R$ is \emph{antisymmetric} if relatedness in both directions
        forces equality:
        \[
          \forall a,b \in A,\quad
          a \mathrel{R} b \;\land\; b \mathrel{R} a \;\Rightarrow\; a = b.
        \]
  \item $R$ is \emph{asymmetric} if relatedness in one direction forbids
        it in the other:
        \[
          \forall a,b \in A,\quad
          a \mathrel{R} b \;\Rightarrow\; \lnot(b \mathrel{R} a).
        \]
\end{enumerate}
\end{definitionbox}

\begin{remark}[Fully quantified form]
Symmetric: $(\forall a,b \in A)\;[a \mathrel{R} b \Rightarrow b \mathrel{R} a]$.
Antisymmetric: $(\forall a,b \in A)\;
[a \mathrel{R} b \land b \mathrel{R} a \Rightarrow a = b]$.
Asymmetric: $(\forall a,b \in A)\;
[a \mathrel{R} b \Rightarrow \lnot(b \mathrel{R} a)]$.
\end{remark}

\begin{remark}[Asymmetric implies irreflexive]
Every asymmetric relation is irreflexive.
If $a \mathrel{R} a$ held, asymmetry would require $\lnot(a \mathrel{R} a)$ —
a contradiction.
Strict inequality $<$ on $\mathbb{Z}$ is asymmetric (hence irreflexive);
$\leq$ on $\mathbb{Z}$ is antisymmetric but not asymmetric.
\end{remark}

\begin{remark}[Symmetric vs.\ antisymmetric]
These are independent.
Equality $=$ is both symmetric and antisymmetric.
The universal relation $A \times A$ (for $|A| > 1$) is symmetric but
not antisymmetric.
Divisibility $\mid$ on $\mathbb{N}$ is antisymmetric but not symmetric.
\end{remark}

\begin{definitionbox}{Definition (Transitive)}
\label{def:rel-transitive}
$R$ is \emph{transitive} if related pairs can be chained:
\[
  \forall a,b,c \in A,\quad
  a \mathrel{R} b \;\land\; b \mathrel{R} c \;\Rightarrow\; a \mathrel{R} c.
\]
\end{definitionbox}

\begin{remark}[Fully quantified form]
$(\forall a,b,c \in A)\;
[a \mathrel{R} b \land b \mathrel{R} c \Rightarrow a \mathrel{R} c]$.
\end{remark}

\begin{remark}[Intuition]
Transitivity says the relation has no ``gaps'' in chains.
It is the property that makes $\leq$ and $\mid$ well-behaved for reasoning.
Friendship is the canonical example of a symmetric, non-transitive relation.
\end{remark}

\begin{definitionbox}{Definition (Total / Connected)}
\label{def:rel-total}
$R$ is \emph{total} (also called \emph{connected} or \emph{trichotomous})
if every pair of distinct elements is related in at least one direction:
\[
  \forall a,b \in A,\quad
  a \mathrel{R} b \;\lor\; b \mathrel{R} a.
\]
\end{definitionbox}

\begin{remark}[Fully quantified form]
$(\forall a,b \in A)\;[a \mathrel{R} b \lor b \mathrel{R} a]$.
Note the disjunction includes $a = b$; if $R$ is also reflexive, this
clause is automatically satisfied when $a = b$.
\end{remark}

% ---------------------------------------------------------
% Composition
% ---------------------------------------------------------

\begin{definitionbox}{Definition (Composition of Relations)}
\label{def:rel-composition}
Let $R \subseteq A \times B$ and $S \subseteq B \times C$.
The \emph{composition} $S \circ R \subseteq A \times C$ is defined by
\[
  a \;(S \circ R)\; c
  \;\iff\;
  \exists b \in B,\quad a \mathrel{R} b \;\land\; b \mathrel{S} c.
\]
\end{definitionbox}

\begin{remark}[Fully quantified form]
$(\forall a \in A)(\forall c \in C)\;
[a\,(S\circ R)\,c
\;\iff\;
(\exists b \in B)\;(a \mathrel{R} b \land b \mathrel{S} c)]$.
\end{remark}

\begin{remark}[Composition is an operation on relations]
When $R$ and $S$ are both on the same set $A$ (i.e.\ $A = B = C$),
composition $\circ$ is a binary operation on the set of all relations on $A$.
It is associative but not in general commutative.
This is the first instance of a binary operation on a \emph{collection of
relations} rather than on elements of a set — a pattern that recurs in
function composition and later in group theory.
\end{remark}

% ---------------------------------------------------------
% Structural classes
% ---------------------------------------------------------

\begin{definitionbox}{Definition (Equivalence Relation)}
\label{def:equiv-relation}
A binary relation $R$ on $A$ is an \emph{equivalence relation} if it is
\begin{enumerate}[label=(\roman*)]
  \item reflexive: $\forall a \in A,\quad a \mathrel{R} a$,
  \item symmetric: $a \mathrel{R} b \Rightarrow b \mathrel{R} a$,
  \item transitive: $a \mathrel{R} b \land b \mathrel{R} c \Rightarrow a \mathrel{R} c$.
\end{enumerate}
We write $a \sim b$ when $R$ is an equivalence relation and the
particular relation is understood from context.
\end{definitionbox}

\begin{remark}[Fully quantified form]
$R$ is an equivalence relation on $A$
$\iff$
$[(\forall a)\, a \mathrel{R} a]$
$\land\;
[(\forall a,b)\, a \mathrel{R} b \Rightarrow b \mathrel{R} a]$
$\land\;
[(\forall a,b,c)\, a \mathrel{R} b \land b \mathrel{R} c \Rightarrow a \mathrel{R} c]$.
\end{remark}

\begin{remark}[Algebraic role]
An equivalence relation on $A$ is precisely the data needed to partition
$A$ into disjoint, exhaustive classes — the equivalence classes.
This partition structure is the algebraic mechanism behind
quotient constructions throughout algebra:
quotient groups, quotient rings, quotient vector spaces, and quotient topologies
are all built by choosing an equivalence relation compatible with the relevant
operations.
\end{remark}

\begin{definitionbox}{Definition (Partial Order)}
\label{def:partial-order}
A binary relation $R$ on $A$ is a \emph{partial order} if it is
\begin{enumerate}[label=(\roman*)]
  \item reflexive: $\forall a \in A,\quad a \mathrel{R} a$,
  \item antisymmetric: $a \mathrel{R} b \land b \mathrel{R} a \Rightarrow a = b$,
  \item transitive: $a \mathrel{R} b \land b \mathrel{R} c \Rightarrow a \mathrel{R} c$.
\end{enumerate}
A set $A$ equipped with a partial order $R$ is a \emph{partially ordered set}
(or \emph{poset}), written $(A, R)$.
We typically write $\leq$ for a partial order.
\end{definitionbox}

\begin{remark}[Fully quantified form]
$(A, \leq)$ is a poset
$\iff$
$[(\forall a)\, a \leq a]$
$\land\; [(\forall a,b)\, a \leq b \land b \leq a \Rightarrow a = b]$
$\land\; [(\forall a,b,c)\, a \leq b \land b \leq c \Rightarrow a \leq c]$.
\end{remark}

\begin{remark}[Strict partial order]
\label{def:strict-partial-order}
The \emph{strict} version $<$ associated to a partial order $\leq$ is defined by
\[
  a < b \;\iff\; a \leq b \;\land\; a \neq b.
\]
The strict relation $<$ is irreflexive, asymmetric, and transitive —
equivalently, it is a \emph{strict partial order}.
Every partial order determines a strict partial order and vice versa:
the two presentations are interconvertible and carry the same information.
\end{remark}

\begin{definitionbox}{Definition (Total Order)}
\label{def:total-order}
A partial order $\leq$ on $A$ is a \emph{total order} (also
\emph{linear order}) if it is also total:
\[
  \forall a,b \in A,\quad a \leq b \;\lor\; b \leq a.
\]
A set $A$ equipped with a total order is a \emph{totally ordered set}
(or \emph{chain}).
\end{definitionbox}

\begin{remark}[Fully quantified form]
$(A, \leq)$ is a chain
$\iff$
$(A, \leq)$ is a poset
$\land\; (\forall a,b \in A)\,[a \leq b \lor b \leq a]$.
\end{remark}

\begin{remark}[Compatibility with operations]
A relation $R$ on $A$ is \emph{compatible with a binary operation}
$\star : A \times A \to A$ if
\[
  a_1 \mathrel{R} a_2 \;\land\; b_1 \mathrel{R} b_2
  \;\Rightarrow\;
  (a_1 \star b_1) \mathrel{R} (a_2 \star b_2).
\]
This condition connects the relational and operational layers of algebraic
structure.
An equivalence relation satisfying it is called a \emph{congruence relation}.
An order satisfying it (for $+$ on an ordered field) yields the ordered field
axioms.
\end{remark}
